\section{長時間エージェント実行のユーザーデモ}

Grok 4.6などを用いて、単発のプロンプトではなく長時間にわたる作業を継続させるデモが共有された。機械的な作業やゲーム構築といった例を通じて、利用者が途中でsteeringしながらエージェントへまとまった仕事を任せる運用が具体化している。

ここで評価されているのは、最高点のベンチマークだけではない。必要な介入量、トークン効率、実行コスト、途中で破綻せず作業を続けられるかといった運用上の性質が重要になっている。エージェント利用がデモから日常的な作業単位へ移る過程を示す投稿群といえる。
\dailyxsource{RobitOverload (@10\_X\_eng)}{https://x.com/10_X_eng/status/2088772041772044556}
